% VI35-DETAILED-PROBLEM-20
\begin{problem}[VI.35.P20: The maximal domain of a graded-map morphism]\label{prob:vi35-20}
Let \(\varphi:S\to T\) be degree preserving. Define
\[
U=\Proj T\setminus V_+\!\bigl(\varphi(S_+)T\bigr).
\]
Prove that contraction and the degree-zero chart maps glue to a morphism
\[
U\longrightarrow\Proj S,
\]
and that \(U\) is maximal for this contraction construction.

\textbf{Provenance.} Canonical VI/35 extension.
\end{problem}

\ifdefined\FullProblemDossiers
\begin{solution}
A point \(\mathfrak q\in\Proj T\) belongs to \(U\) exactly when there exists homogeneous
\[
f\in S_+
\]
with
\[
\varphi(f)\notin\mathfrak q.
\]
Then
\[
\varphi^{-1}(\mathfrak q)
\]
is a homogeneous prime of \(S\) missing \(f\), hence a point of \(D_+(f)\subseteq\Proj S\).

Accordingly,
\[
U
=
\bigcup_{f\in S_+\text{ homogeneous}}D_+(\varphi(f)).
\]
On the open \(D_+(\varphi(f))\), the graded map induces a degree-zero ring map
\[
S_{(f)}
\longrightarrow
T_{(\varphi(f))},
\qquad
\frac a{f^m}\longmapsto
\frac{\varphi(a)}{\varphi(f)^m}.
\]
Contravariance of \(\Spec\) gives a morphism
\[
D_+(\varphi(f))
\longrightarrow
D_+(f).
\]
On overlaps these maps agree because they are induced by the same original homomorphism after localization in
stages. Hence they glue to
\[
\boxed{U\to\Proj S}.
\]

If \(\mathfrak q\notin U\), then
\[
\varphi(S_+)T\subseteq\mathfrak q,
\]
so
\[
S_+\subseteq\varphi^{-1}(\mathfrak q).
\]
The contracted prime is therefore not a point of \(\Proj S\). No morphism defined by contraction can include this
point. Thus \(U\) is maximal for this construction.
\end{solution}
\fi
